\section{Evaluation protocol}
\label{sec:protocol}

Fish has a single chance node --- the deal --- followed by a public transcript,
and that structure permits an unusually strong variance-reduction design. This
section fixes the experimental unit, the interval construction, the
fitting/evaluation separation, and the metrics of \S\ref{sec:results}.

It also fixes something the two preceding studies did not, and the omission cost
this one five mechanisms. A protocol section is normally read as bookkeeping. In
this study it is a result: \S\ref{sec:protocol-resolution} shows that the
budget most of the corpus's published ablation rows were run at cannot separate
an effect of the size those rows report from a seed draw, and that five separate
mechanisms cleared a nominal \vsixBootAlpha{}\% paired interval at that budget and returned
exactly zero when re-run larger on two disjoint banks. Every claim in this report
is paired and re-run on a second bank for that reason.

A reader should be able to tell, for every number in \S\ref{sec:results}, whether
it is paired or unpaired, which bank it came from, and whether that bank was used
in fitting. The earlier reports could not support that reading, which is why
several of their comparisons are corrected in \S\ref{sec:corrections}.

\subsection{Six-rotation duplicate blocks}
\label{sec:protocol-blocks}

Seats alternate between the two teams, so the cyclic group of six seat rotations
of a fixed deal forms a balanced block: three seat shifts by two team
orientations. The odd rotations exchange which team holds which hand-triple; the
even rotations permute hands within a team. Playing all \vsixRotations{}
rotations of every deal therefore has each policy hold every hand-triple of that
deal exactly once, so the intrinsic advantage of the deal cancels from the
comparison rather than merely being averaged down. A matchup reported over $D$
deals is $6D$ games organised into $D$ blocks of six. This is the duplicate-block
convention of competitive play \cite{acpc}, strengthened: the competition design
balances the team label, and the six-seat rotation additionally balances the
assignment of individual hands within a team.

One guarantee of that convention does not survive, and it is stated here rather
than in a footnote because it silently halves several denominators. In a
\emph{mirror} match --- both teams running the same policy specification --- the
two team orientations of a deal produce byte-identical games, so each game is
played twice and the effective sample is \vsixMirrorDuplication{}\% of the
nominal one. Rates are unaffected; denominators are not. Every table in
\S\ref{sec:results} states which convention its $n$ uses, and a mirror
\texttt{match} at two rotations returns exactly $50\%$ for that reason and not
because the policies are balanced.

\subsection{Deal-clustered bootstrap}
\label{sec:protocol-ci}

The six rotations of one deal are one correlated cluster: they share the same
card distribution and differ only in seat assignment. Resampling \emph{games}
would treat those six observations as independent and understate the variance.
Every interval in this report therefore resamples \emph{deals} --- a cluster
bootstrap of the per-deal win count for a single arm, and a paired bootstrap of
the per-deal difference for a comparison between two configurations played on the
same deals --- taking the $2.5\%$ and $97.5\%$ quantiles of \vsixBootReps{}
resamples, for \vsixBootAlpha{}\% intervals. Both are implemented in
\filepath{engine/src/arena.hpp}.

These are \emph{percentile} bootstrap intervals with the deal as the resampling
cluster, and they inherit the known coverage error of that construction. We keep
it for two reasons. Clustering by deal is the correction that matters most in
this design, because the within-deal correlation across the six rotations is
large and ignoring it understates the variance by a substantial factor, whereas
the percentile-versus-BCa distinction is second-order at these sample sizes. And
we have not implemented the bias-corrected and accelerated or bootstrap-$t$
corrections, so no claim of nominal coverage is made: intervals should be read as
approximate, and a difference whose interval sits close to zero should not be
treated as resolved. \S\ref{sec:protocol-resolution} is what that caution looks
like when it is applied rather than merely stated.

\subsection{Paired estimators}
\label{sec:protocol-paired}

The headline comparison is a paired difference on common random numbers, and
every ablation is paired against its own control. The pairing is what makes the
design competitive with the estimator-side variance reduction developed for
poker \cite{zinkevich-divat,burch-aivat} without needing an auxiliary value
function: in this game the deal is the whole chance node, so replaying it cancels
the dominant variance component directly. \S\ref{sec:fit-paired} gives the
measured justification, and the same estimator is used in fitting and in
evaluation so that a candidate is scored during the search by the quantity it
will be judged on afterwards.

One quantity in this report is not available paired. Throughput
(\S\ref{sec:results-throughput}) is a property of one binary on one machine
rather than of a matchup, so there is no common deal to pair on; it is reported
as a single-arm figure measured under load and is a lower bound for that reason.

\subsection{The resolution requirement}
\label{sec:protocol-resolution}

Five mechanisms in this study cleared a nominal \vsixBootAlpha{}\% paired
interval at \numResSmallBudgetLo{}--\numResSmallBudgetHi{} games per cell and
returned exactly zero at
\numResLargeBudget{} games per cell on two disjoint banks.
Table~\ref{tab:resolution} is the record.

\begin{table}[ht]
\centering
\caption{Five mechanisms that resolved at a small budget and did not replicate.
Sign convention is the corpus's: $\delta = \text{reference} - \text{variant}$, so
a \emph{negative} $\delta$ means the variant is better. Every entry is a paired
per-deal margin with a deal-clustered percentile interval. The large-budget
column gives two independent seed banks. Source:
\texttt{research/v06/notes/R12-mechanism-trials.md}, \S6.1.}
\label{tab:resolution}
\footnotesize
\begin{tabular}{@{}>{\raggedright\arraybackslash}p{0.235\linewidth}>{\raggedright\arraybackslash}p{0.265\linewidth}>{\raggedright\arraybackslash}p{0.42\linewidth}@{}}
\toprule
Mechanism & Small-budget paired $\delta$ & Large-budget paired $\delta$, two banks \\
\midrule
\texttt{wVoid} at $1.5$
  & \numResVoidSmall{} [\numResVoidSmallCI{}], $n{=}\numResVoidSmallN{}$
  & \numResVoidLargeA{} [\numResVoidLargeACI{}] and
    \numResVoidLargeB{} [\numResVoidLargeBCI{}], $n{=}\numResVoidLargeN{}$ \\
\addlinespace
\texttt{wVoid} at $3$
  & \numResVoidThreeSmall{} [\numResVoidThreeSmallCI{}], $n{=}\numResVoidThreeSmallN{}$
  & as above \\
\addlinespace
delete the perseveration bonus
  & \numResPerseverSmall{} [\numResPerseverSmallCI{}], $n{=}\numResPerseverSmallN{}$
  & \numResPerseverLargeA{} [\numResPerseverLargeACI{}] and
    \numResPerseverLargeB{} [\numResPerseverLargeBCI{}], $n{=}\numResPerseverLargeN{}$ \\
\addlinespace
\texttt{declareMargin} more negative
  & \numResDeclSmall{}~pp unpaired, $n{=}\numResDeclSmallN{}$
  & \numResDeclLarge{} [\numResDeclLargeCI{}] --- the reference is better \\
\addlinespace
delete the negative value-of-information term
  & \numResVoiSmall{} [\numResVoiSmallCI{}], $n{=}\numResVoiSmallN{}$
  & not separated \\
\bottomrule
\end{tabular}
\end{table}

The planning rule that follows from this is arithmetic, not a convention. For an
unpaired win-rate comparison near even the half-width of a nominal $95\%$
interval is
\begin{equation}
  1.96\sqrt{\frac{p(1-p)}{N}} \;\approx\; \frac{0.98}{\sqrt{N}}
  \quad\text{at } p \approx \tfrac12 ,
  \label{eq:halfwidth}
\end{equation}
that is, about $98/\sqrt{N}$ percentage points at $N$ games. At $N = \numResSmallBudgetLo{}$ that is a
half-width of nearly five points; at $N = \numResVoidLargeN{}$ it is about
three. An effect
of one to two points --- the size every mechanism in Table~\ref{tab:resolution}
appeared to have --- is therefore indistinguishable from a seed draw at the
budget most of this corpus's published ablation rows were run at. Pairing buys
back a large factor, but it buys it back from the deal variance and not from the
requirement to check that the result is not the bank.

Two rules follow and both are enforced in this report. Every claim is
\emph{paired}: the variant and its control are played on the same deals and the
interval is on the per-deal difference. And every claim is \emph{re-run on a
second, disjoint bank}, and reported as replicated only if it holds in sign and
size on both. Where a result resolves on one bank and not the other, that is
said, and the per-cell rows are given so a reader can see which half of the
comparison carries it.

The general lesson is uncomfortable and is stated plainly in
\S\ref{sec:limitations}: at the budget this project used before v0.6, a
mechanism that does nothing has a substantial chance of being published as a
mechanism that does something, and several were.

\subsection{Seed separation and what was held out}
\label{sec:protocol-seeds}

Fitting used base seeds \vsixFitASeed{} (run~A) and \vsixFitCSeed{} (run~C).
Every seed bank used from \S\ref{sec:results} onwards is disjoint from both by
construction, and the disjointness is enforced by a registry of the
\vsixReservedSeeds{} reserved seeds that the battery checks, rather than by
discipline. The evaluation banks are \texttt{90210}, \texttt{31337},
\texttt{515151}, \texttt{777001} and \texttt{424242} for the head-to-head;
\texttt{515253} for the per-style panel; \texttt{606060} for the frozen-policy
ablations; \texttt{717171} for calibration; \texttt{828282} for the rule
dialects; \texttt{313131} for the partner regimes; \texttt{90210} again for the
search ladder, the deviation-rate probe and the pre-gate audit; \texttt{31} for
the pathology KPIs and the tie and belief probes; and \texttt{515253} with
\texttt{6543210} for the exploitability probe's fitting and evaluation halves
respectively. \vsixHeldOutBanks{} of them are held out and \vsixFitBanks{} are
fitting banks. Banks carried forward from the v0.5 study are named in
\S\ref{app:repro-commands} so that the corresponding rows are directly
comparable.

The separation is narrower than ``held out'' usually implies, and the report
scopes its claims accordingly. Deals are held out: no deal used in any reported
experiment appears in any fitting bank. Opponent \emph{identities} largely are
not --- the fitting panels of \S\ref{sec:fit-schedule} include \vfive{},
\vthree{} and two scripted archetypes that also appear in the evaluation panel
--- so the head-to-head and per-style results measure generalisation to new deals
rather than to new strategic populations. The exceptions, which are held out in
both senses, are named in \S\ref{sec:results-panel}.

\subsection{Metrics}
\label{sec:protocol-metrics}

The primary metric is win rate, reported with its paired margin over the control.
Sets differential --- mean half-suits won out of nine --- accompanies it, because
a policy can win while converting a smaller fraction of its asks and this report
contains a matchup in which that happens. Ask hit rate and declaration accuracy
are diagnostic rather than primary: they describe the decision quality of one
component and a policy can trade either against the other.

The pathology KPIs of \S\ref{sec:fit-gates} are gates, not metrics. They do not
enter any comparison and no configuration is ranked on them; they decide only
whether a configuration is eligible to be measured at all.

One standing rule governs how the opponent panel is reported: \emph{never
headline an aggregate over it}. The headline is the worst case and the minimax
regret, with the full per-opponent profile printed alongside. A mean over a
scripted panel is a weighted average whose weights are the panel's composition,
and the panel's composition is a choice this study made. For the same reason
rating systems are deliberately not used as a headline here
\cite{coulom-whr,herbrich-trueskill,nelo,balduzzi-reeval}: a single scalar over a
fixed scripted panel is exactly the aggregate the rule forbids, and the
transitivity it assumes is not established for this population. That is stated
rather than the systems being omitted silently.
